Large language models, trained through massive text compression, may discover structural analogies between concept domains that differ from those humans naturally use.
We investigate whether LLMs internally create and leverage such \emph{unexpected partial isomorphisms}---structure-preserving mappings between concepts that are distant in human similarity judgments but close in model representation space.
We conduct three complementary experiments.
First, we compare pairwise concept similarities from LLM embeddings against WordNet-derived human similarity scores across 6,600 cross-domain concept pairs, identifying 162 significantly surprising pairs (exceeding $2\sigma$ above the mean surprise score).
Second, we test whether these unexpected analogies improve reasoning: when provided as hints to \gptfour, they yield higher-quality insights (3.62/5) than conventional analogies (3.38/5) or no hints (3.00/5).
Third, we analyze MLP activations in \pythia and find that surprising concept pairs share significantly more activation overlap than matched controls (Cohen's $d = 0.81$, $p = 0.013$), with the effect concentrated in semantic processing layers 5--17.
Our results show that LLMs have discovered a rich space of ``alien analogies''---functional and structural parallels like \emph{ingredient}$\leftrightarrow$\emph{enzyme} and \emph{foundation}$\leftrightarrow$\emph{anchor}---that humans do not commonly use but that reflect genuine shared structure.
These findings suggest that neural compression creates emergent conceptual organization that could yield novel metaphors for human reasoning.
